\documentclass[aspectratio=169,11pt]{beamer}

\usetheme{Madrid}
\usecolortheme{default}
\usefonttheme{professionalfonts}
\setbeamertemplate{navigation symbols}{}
\setbeamertemplate{footline}[frame number]

\usepackage{amsmath, amssymb, booktabs}

\newcommand{\E}{\mathbb{E}}
\newcommand{\1}{\mathbf{1}}
\newcommand{\Prb}{\mathbb{P}}

\title{Curating Network Data And Descriptive Statistics}
\subtitle{Empirical Methods --- Week 18}
\author{Teaching deck draft}
\date{}

\begin{document}

\begin{frame}
  \titlepage
\end{frame}

\begin{frame}{Week 18 question}
\begin{itemize}
    \item How should researchers define, measure, and summarize relationships that shape economic outcomes?
    \item Network curation comes before network causal inference.
    \item Node, edge, weight, direction, timing, and boundary choices shape the estimand.
    \item A graph is not evidence until the link meaning is justified and auditable.
\end{itemize}

\medskip
This is the opening lecture of the optional network-methods block.
\end{frame}

\begin{frame}{Nodes, edges, and bipartite data}
\small
\[
A_{ij}=
\begin{cases}
1 & \text{if there is a link from } i \text{ to } j,\\
0 & \text{otherwise.}
\end{cases}
\qquad
B_{if}=
\begin{cases}
1 & \text{if worker } i \text{ is linked to firm } f,\\
0 & \text{otherwise.}
\end{cases}
\]

\[
P^{worker}_{ij}=\sum_f B_{if}B_{jf}
\]

\begin{itemize}
    \item Nodes may be workers, firms, schools, managers, households, vacancies, or neighborhoods.
    \item Edges may be referrals, coworking, co-residence, communication, assignment, or transactions.
    \item Direction and weight matter when recommendations, supervision, duration, or intensity are asymmetric.
    \item Bipartite data are often the raw object; projection creates a new empirical object.
\end{itemize}
\end{frame}

\begin{frame}{Network definitions are substantive}
\small
\begin{tabular}{p{0.22\linewidth} p{0.32\linewidth} p{0.34\linewidth}}
\toprule
Choice & Empirical version & Economic content \\
\midrule
Node & worker, firm, vacancy, household & whose choices and outcomes matter \\
Edge & referral, coworker, neighbor, message & information, trust, exposure, monitoring \\
Direction & \(i \to j\), \(j \to i\), undirected & asymmetric recommendation or influence \\
Weight & duration, frequency, value & strength of relationship or exposure \\
Timing & month, spell, year, cohort & whether exposure precedes outcome \\
Boundary & firm, city, platform, survey frame & which links are unobserved or truncated \\
\bottomrule
\end{tabular}

\medskip
The same raw records can imply different network variables under different mechanisms.
\end{frame}

\begin{frame}{Labor economics uses}
\begin{itemize}
    \item Referrals: who hears about vacancies and who is recommended.
    \item Information transmission: how job knowledge travels through contacts, neighborhoods, and coworkers.
    \item Coworker exposure: mentoring, norms, monitoring, team learning, and job ladders.
    \item Neighborhood job networks: local residents connected to employers and opportunities.
    \item Learning and diffusion: practices, technologies, beliefs, and search strategies.
    \item Bargaining and inequality: unequal access to high-value contacts and outside options.
\end{itemize}

\medskip
Anchor papers: Autor; Bayer, Ross, and Topa; Beaman and Magruder; Barwick et al.; Hellerstein et al.
\end{frame}

\begin{frame}{Data sources and construction}
\small
\begin{tabular}{p{0.28\linewidth} p{0.30\linewidth} p{0.30\linewidth}}
\toprule
Source & Natural graph & Main limitation \\
\midrule
Linked employer--employee & worker--firm bipartite; coworker projection & informal ties outside firms \\
Referral-program records & directed employee--applicant links & only formal referrals \\
Survey rosters & friendship, kinship, advice, support & recall error and missing weak ties \\
Platform logs & worker--client or interaction graph & platform rules and proprietary access \\
School or cohort rosters & peer exposure graph & assignment and co-presence are not random \\
Neighborhood records & local opportunity or co-residence graph & proximity is not interaction \\
\bottomrule
\end{tabular}

\medskip
Preserve the raw relational object before computing projected links or summaries.
\end{frame}

\begin{frame}{Descriptive statistics are theory-laden}
\small
\[
d_i=\sum_j A_{ij}
\qquad
s_i=\sum_j W_{ij}
\qquad
c_i=\lambda^{-1}\sum_j A_{ij}c_j
\]

\[
H=
\frac{\sum_i\sum_j A_{ij}\mathbf{1}\{g_i=g_j\}}
{\sum_i\sum_j A_{ij}}
\]

\begin{itemize}
    \item Degree can mean opportunity access, referral reach, survey visibility, or firm size.
    \item Strength adds intensity but can hide whether exposure is broad or concentrated.
    \item Centrality can proxy information access only when paths carry the relevant information.
    \item Homophily describes similarity among linked nodes; it does not separate sorting from influence.
\end{itemize}
\end{frame}

\begin{frame}{Pitfalls and diagnostics}
\small
\begin{tabular}{p{0.27\linewidth} p{0.60\linewidth}}
\toprule
Risk & Diagnostic \\
\midrule
Missing links & distinguish unobserved ties from true zero ties \\
Endogenous formation & state whether the network is an outcome, exposure, or design object \\
Boundary truncation & test alternative graph frames and inspect isolates/components \\
Projection error & compare bipartite statistics with projected network statistics \\
Homophily vs influence & avoid causal language until formation and interference are handled \\
Privacy & report how masking changes degree, strength, and components \\
Dynamic mismatch & vary time windows and check exposure precedes outcome \\
\bottomrule
\end{tabular}
\end{frame}

\begin{frame}{Practical rules}
\begin{enumerate}
    \item State the labor-market mechanism before defining the graph.
    \item Define node, edge, direction, weight, layer, timing, and boundary separately.
    \item Preserve raw edge lists, bipartite spell tables, and survey rosters.
    \item Build the primary network statistic and one plausible alternative.
    \item Report missing-link risk and privacy masking as measurement issues.
    \item Normalize projections when institution size mechanically creates links.
    \item Test whether patterns are driven by hubs, components, or boundary truncation.
    \item Keep descriptive diagnostics separate from causal robustness.
\end{enumerate}
\end{frame}

\begin{frame}{Applied paper anchors}
\begin{itemize}
    \item Autor: labor markets are wired through search channels and intermediaries.
    \item Bayer, Ross, Topa: residential proximity can reveal informal hiring networks.
    \item Beaman, Magruder: referral incentives affect worker selection.
    \item Friebel et al.: employee referral programs as a management practice.
    \item Barwick et al.: referrals can shape labor-market inequality.
    \item Hellerstein, McInerney, Neumark: networks matter after mass layoffs.
\end{itemize}

\medskip
The lesson is not ``compute centrality.'' The lesson is: define the economic relation.
\end{frame}

\begin{frame}{Research Lab design}
\begin{columns}[T,onlytextwidth]
\begin{column}{0.32\linewidth}
\textbf{Reproduce}
\begin{itemize}
    \item synthetic workers
    \item neighborhoods
    \item employers
    \item bipartite links
    \item local exposure
\end{itemize}
\end{column}
\begin{column}{0.32\linewidth}
\textbf{Diagnose}
\begin{itemize}
    \item boundary choices
    \item edge rules
    \item binary vs weighted
    \item missing links
    \item projection risk
\end{itemize}
\end{column}
\begin{column}{0.32\linewidth}
\textbf{Transfer}
\begin{itemize}
    \item directed referrals
    \item in-degree/out-degree
    \item referral success
    \item inequality memo
    \item causal preview
\end{itemize}
\end{column}
\end{columns}
\end{frame}

\begin{frame}{Where curation stops}
\begin{itemize}
    \item Curation defines edge lists, adjacency matrices, bipartite objects, and summaries.
    \item Curation can diagnose measurement sensitivity, missing links, and graph boundaries.
    \item It does not solve reflection, sorting, peer effects, or interference.
    \item Network causal inference begins when the researcher claims counterfactual effects through links.
    \item Lecture 19 takes up exposure mappings, peer effects, network experiments, and dependence.
\end{itemize}
\end{frame}

\begin{frame}{Takeaways}
\begin{itemize}
    \item Network definition is part of the estimand.
    \item A link is an economic channel, not just a row in an edge list.
    \item Bipartite and multilayer objects should not be collapsed without a reason.
    \item Degree, centrality, clustering, and homophily are descriptive until the mechanism is defended.
    \item Missing links, boundaries, privacy, and timing can change the object.
    \item Strong curation makes Lecture 19's causal designs possible.
\end{itemize}
\end{frame}

\end{document}
